\chapter{Evaluation and Results}\label{ch:eval}

\section*{Introduction}
\addcontentsline{toc}{section}{Introduction}

Everything in this report so far has been a claim about design: here is how the
loop works, here is how it grounds itself, here is what it refuses to do. None
of that answers the question anyone actually cares about, which is whether the
thing works, and how often. This chapter answers that with numbers, including
the ones that are not flattering.

\section{Method}

Measuring an agent like this is harder than measuring a classifier, because
there is no single correct string to compare against. What there \emph{is}, for
an assistant sitting on top of a database, is a correct \emph{number}: ask how
many subscribers are on 5G and the answer either contains the value the database
holds or it does not. That is what makes an objective golden set possible
without anybody having to grade answers by hand.

The harness holds sixteen questions across seven categories, and each one is
paired with a programmatic check rather than an expected string:

\begin{itemize}
    \item \textbf{Metric} --- an objective count has to appear in the answer.
    \item \textbf{Categorical} --- an objective label has to appear.
    \item \textbf{Distribution} --- the answer has to carry a chart.
    \item \textbf{Hallucination} --- the question asks for a metric the schema
    does not contain, so a correct answer says so instead of producing a number.
    \item \textbf{Safety} --- destructive instructions and prompt injections
    must not be executed.
    \item \textbf{Routing} --- asking to \emph{create} something has to produce
    a proposal for human approval, not a silent action.
    \item \textbf{Empty} --- an impossible filter has to degrade gracefully
    rather than crash.
\end{itemize}

Reference values get computed straight from the databases immediately before
each question is asked, and again immediately after, and an answer counts as
correct if it lands anywhere in that window. This is not a tolerance I invented
to make the numbers look better: the simulator never stops, so the true value
genuinely moves during the few seconds an answer takes to produce, and that
window is the real range it occupied. Memory is reset between cases too, so no
question gets to benefit from the one before it.

\section{Results}

The harness was run twice, a couple of days apart, with no change to the agent
in between. Both runs are reported here, because the difference between them
turned out to be as informative as either one on its own.

\begin{table}[H]
\centering
\begin{tabular}{lccl}
\toprule
\textbf{Category} & \textbf{Run 1} & \textbf{Run 2} & \textbf{Interpretation} \\
\midrule
Metric        & 5/5 & 5/5 & Factual retrieval is reliable \\
Categorical   & 1/1 & 1/1 & \\
Safety        & 2/2 & 2/2 & Refused deletion and prompt injection \\
Empty         & 2/2 & 2/2 & No crash on impossible filters \\
Distribution  & 0/1 & 0/1 & Asked for clarification instead \\
Routing       & 0/1 & 0/1 & Analysed instead of proposing \\
Hallucination & 1/4 & 2/4 & \textbf{The weak point} \\
\midrule
\textbf{Total} & \textbf{11/16} & \textbf{12/16} & \textbf{68.75\% and 75\%} \\
\bottomrule
\end{tabular}
\caption{Golden-set results by category, across two runs of the same sixteen
questions against the same agent.}
\label{tab:eval-results}
\end{table}

The shape of that table matters more than the total. Everything the agent was
built to do first --- pull back a fact, refuse something dangerous, survive a
question with no answer --- it does without a single miss. What it does not do
reliably is decline to answer when the data cannot support an answer.

\section{Failure Analysis}

Four or five failures depending on the run, and they are not the same kind of
failure at all. What follows describes each of them; the counts above say how
often each one showed up.

\subsection{Swapping In a Neighbouring Metric}

Three of the five failures in run 1 come from one behaviour, and two of them
recur in run 2. Asked for the \emph{call setup
success rate}, which no column in either database holds, the agent answered
98.89\%, worked out from the average dropped-call rate. Asked for \emph{customer
lifetime value by segment}, it came back with average ARPU by segment. Asked how
many subscribers are on \emph{fibre}, it answered with fixed-wireless numbers.

None of these is a fabrication in the usual sense. Each one is a real figure,
correctly computed, from a real column --- answering a question nobody asked.
This is the limitation from the conclusion made concrete: the verification layer
checks that every number in an answer appears in the query results, and in all
three cases it did. The layer has no way of knowing the question wanted
something else entirely.

Worth pointing out what it got right in the same category, though: asked for the
\emph{6G adoption rate}, it said there is no 6G in the network and that 5G is
the highest technology present. The difference is that there is no plausible
neighbouring column for 6G, so there was nothing to swap in.

\subsection{The Same Question, Three Different Inventions}

The clearest single result in this whole evaluation came from running it more
than once. Asked for the \emph{call setup success rate} --- a metric no column
in either database holds --- the agent answered differently every single time:

\begin{table}[H]
\centering
\begin{tabular}{ll}
\toprule
\textbf{Run} & \textbf{Answer given} \\
\midrule
1 & The call setup success rate is 38.77\%, so nearly 61\% of attempts are failing \\
2 & The call setup success rate is 98.89\%, based on the average dropped call rate \\
3 & The call setup success rate is 0\%, which indicates a critical issue \\
\bottomrule
\end{tabular}
\caption{Three runs, one question, three fabricated values for a metric that
does not exist.}
\label{tab:callsetup-variance}
\end{table}

Nothing changed between those runs. Not the agent, not the prompt, not the
schema. All three answers are confident, none of them is real, and the third one
escalates a number it invented into a critical issue somebody should look at.

This is the non-determinism limitation made concrete. It also shows why a single
pass through the golden set is an estimate rather than a measurement: the fibre
question failed in run 1, where the agent answered with fixed-wireless numbers,
and passed in run 2, where it correctly said there are no subscribers on fibre.
Same question, same agent, opposite outcome --- which is the one point in this
chapter I would want a reader to take away over the totals.

\subsection{Clarification Counted as a Failure}

Asked to show the subscriber distribution by technology, the agent came back
with a clarification prompt instead of a chart. Judged against the test that is
a failure. Judged against the design it is the disambiguation gate from
Chapter~\ref{ch:agent} doing precisely what it was built to do on a breakdown
request. It is recorded as a failure here rather than quietly excused, but it is
a disagreement between the test and the design, not a defect.

\subsection{Routing}

Asked to create a 5G upsell campaign, the agent produced an analysis of the
opportunity instead of a proposal for approval. Nothing got executed --- the
human-in-the-loop guarantee held --- but the request should have been routed to
the approval path and it was not. This is a real gap, and the only one of the
five with an obvious fix.

\section{What Building the Harness Found}

The harness was written to measure the agent. It spent most of its first few
runs finding faults everywhere except the agent, which is itself a decent
argument for building one.

Its very first run failed all sixteen cases with \texttt{Unable to locate
credentials}. The harness had never loaded the environment file that holds them,
because the web server was the only other entry point and it loads them itself.
That one missing line is the whole reason the evaluation in this project had
never been run before.

Two more cases then failed on counts. The agent reported 3{,}162 subscribers on
5G where the reference said 3{,}204, and the agent was right. The reference
counted rows in \texttt{subscriber\_technology} directly; the agent joined to
\texttt{subscribers} the way it should. The gap is 371 orphaned rows across that
table --- technology records whose subscriber no longer exists, left behind by
an older identifier migration that never cascaded. So the evaluation found a
defect in the data, not in the answer.

A third run turned up a crash: a single arrow character in a model response
triggered an encoding failure on Windows, silently dropping the agent from its
tool-calling path down to the older text-parsing fallback. In normal use this is
completely invisible --- the answer still arrives, just by a worse route.

Not one of those three would have been found by asking the assistant questions
and reading the answers, which is how everything else in this project got
verified.

\section{Churn Scoring: A Transfer Failure and What Replaced It}

The churn model is a component of this project rather than the point of it, but
the agent reports its output to users, so whether it is valid matters. Actually
testing that turned out to produce the most useful negative result of the whole
project.

\subsection{The Model, and Its Performance on Its Own Data}

The original model is an XGBoost classifier~\cite{xgboost} trained on a public
Indian telecom dataset~\cite{telecom_churn_dataset} of 99{,}999 subscriber records, of which 9.1\% churned,
using fifteen engineered features dominated by relative decline --- the drop in a
subscriber's spend against their own baseline rather than its absolute level. The
split is 80/20 stratified with SMOTE~\cite{smote} applied to the training half only.

\begin{table}[H]
\centering
\begin{tabular}{lcccc}
\toprule
& \textbf{Precision} & \textbf{Recall} & \textbf{F1} & \textbf{Support} \\
\midrule
Active & 0.96 & 0.97 & 0.96 & 18{,}186 \\
Churn  & 0.62 & 0.56 & 0.59 & 1{,}814 \\
\midrule
\multicolumn{5}{l}{Accuracy 0.9295 \quad AUC-ROC 0.8845 \quad 5-fold CV 0.9255 $\pm$ 0.0413} \\
\bottomrule
\end{tabular}
\caption{XGBoost churn model on the held-out split of its own training dataset.}
\label{tab:churn-perf}
\end{table}

That headline accuracy is the least useful number in the table. At a 9.1\% base
rate, simply predicting that nobody ever churns already scores 0.91. The
AUC-ROC of 0.88 is the honest one: it says the model \emph{ranks} subscribers
well, which is exactly what a percentile-banded system needs it to do.

Recall on the churn class needs the same care. The 0.56 in the table is measured
at the default 0.5 cut, which is not the operating point the system actually
uses --- it bands by percentile instead. Measured at the thresholds that are
really deployed, things look quite different:

\begin{table}[H]
\centering
\begin{tabular}{lcccc}
\toprule
\textbf{Operating point} & \textbf{Threshold} & \textbf{Flagged} & \textbf{Precision} & \textbf{Recall} \\
\midrule
Default binary cut          & 0.500 & 2{,}153 & 0.52 & 0.62 \\
High band (top decile)      & 0.526 & 2{,}051 & 0.54 & 0.61 \\
Medium and above (top 30\%) & 0.145 & 6{,}001 & 0.25 & \textbf{0.84} \\
\bottomrule
\end{tabular}
\caption{The same model at different operating points. Reporting only the default
cut understates what the deployed configuration achieves.}
\label{tab:churn-thresholds}
\end{table}

At the band a retention campaign would genuinely use, the model finds 84\% of
churners rather than 56\%. To keep it simple: a single precision and recall pair
describes a threshold, not a model, and quoting only one of them undersells what
the deployed configuration does.

\subsection{Testing Whether It Transfers}

Everything above is performance on the dataset the model was trained on. Here it
gets applied to a completely different world: the training data is Indian
\emph{prepaid} recharge behaviour, and this environment is \emph{postpaid}
billing. Whether it transfers is an empirical question, so I tested it by
comparing the distribution of every feature as trained against the same feature
as it is actually fed here.

\begin{table}[H]
\centering
\begin{tabular}{lrrr}
\toprule
\textbf{Feature} & \textbf{Training mean} & \textbf{Fed mean} & \textbf{Ratio} \\
\midrule
\texttt{data\_mb\_m8}     & 135.41 & 40{,}807.13 & $301\times$ \\
\texttt{rech\_num\_m8}    &   7.21 &        0.85 & $8.5\times$ \\
\texttt{has\_roaming}     &   0.37 &        0.05 & $7\times$ \\
\texttt{voice\_min\_m8}   & 125.86 &       30.84 & $4\times$ \\
\bottomrule
\end{tabular}
\caption{The four worst feature mismatches. Eight of fifteen features diverge
materially; these are the extremes.}
\label{tab:churn-transfer}
\end{table}

The data-usage feature is the clearest failure of the lot. Its training median
is \textbf{zero} --- most of that prepaid base used no 3G data whatsoever ---
while the median subscriber here gets fed roughly 40~GB. So every tree split
that was learned on that feature is now being evaluated miles outside the range
it was ever fitted on.

Two further defects came out of the same inspection. The mapping sets
\texttt{rech\_amt\_drop\_pct} equal to \texttt{arpu\_drop\_pct} and
\texttt{rech\_amt\_m8} equal to \texttt{arpu\_m8}, so two of the fifteen inputs are
exact duplicates of two others and the model's most important feature is counted
twice. And \texttt{rech\_num\_m8}, a count of recharges in the range 0--30 during
training, is mapped here to a binary paid/unpaid flag.

So the 0.88 AUC is real, and it says nothing at all about this data. The model
was never retrained, revalidated or recalibrated for the environment it got
deployed into. It was just pointed at it and trusted.

\subsection{The Replacement: An Explainable Scorecard}

Rather than keep presenting an uncalibrated model as though it were a
probability, churn scoring got replaced with a transparent scorecard built
entirely from the operator's own columns:

\begin{table}[H]
\centering
\begin{tabular}{lcl}
\toprule
\textbf{Signal} & \textbf{Weight} & \textbf{Source} \\
\midrule
ARPU decline vs own baseline & 0.35 & \texttt{customer\_value.arpu} history \\
Unpaid bills                 & 0.20 & \texttt{billing.payment\_status} \\
Data-usage decline           & 0.15 & \texttt{dou\_monthly} \\
Voice decline                & 0.10 & \texttt{ott\_monthly} \\
Low engagement               & 0.10 & \texttt{dou\_monthly.days\_active} \\
Short tenure                 & 0.10 & \texttt{subscribers.activation\_date} \\
\bottomrule
\end{tabular}
\caption{The churn scorecard. Weights sum to 1.0; the experience penalty of
Chapter~\ref{ch:livingdb} is added on top and capped at 0.40.}
\label{tab:churn-heuristic}
\end{table}

Live service experience --- open incidents and unresolved complaints --- is
deliberately excluded from the scorecard and left to the existing experience
penalty, so that the two are not counted twice. The split is clean: the scorecard
scores commercial behaviour, the penalty scores network experience, and churn risk
is their sum.

Three things make this better than what it replaced. It is computed from data
this operator actually holds, so there is no transfer assumption to worry about
at all. Its bands are percentiles of the live score distribution rather than
constants inherited from somebody else's dataset. And every score breaks down
into a sentence a human can check --- flagged because spend fell against their
own baseline and two bills went unpaid --- which the tree ensemble could never
give us.

\subsection{What the Scorecard Finds, and What It Cannot}

Scored across the base, the bands separate sharply on payment behaviour and only
weakly on anything else:

\begin{table}[H]
\centering
\begin{tabular}{lrrrr}
\toprule
\textbf{Band} & \textbf{Subscribers} & \textbf{Unpaid bills} & \textbf{Avg ARPU} & \textbf{Avg tenure} \\
\midrule
High   & 5{,}036  & 35.3\% & 109.8 & 1{,}406 d \\
Medium & 10{,}051 & 25.2\% & 109.4 & 1{,}434 d \\
Low    & 34{,}958 &  8.7\% & 110.6 & 1{,}481 d \\
\bottomrule
\end{tabular}
\caption{Scorecard bands against independent indicators.}
\label{tab:churn-bands}
\end{table}

Unpaid bills separate by a factor of four across the bands, which is the
scorecard doing its job. Average ARPU does not separate at all, and that is
worth saying out loud rather than hiding: the heaviest term in the whole
scorecard contributes nothing measurable in this environment.

The reason is a property of the simulator, not of the scorecard. Every subscriber's
ARPU is advanced by the same random drift each month, so while the series is
continuous --- month-to-month correlation of 0.998, against 0.15 for the raw
billing amounts, which is why the scorecard reads \texttt{customer\_value} rather
than \texttt{billing} --- there is no subpopulation whose spend systematically
declines. There are no declining customers in this world for a decline term to
find.

That is a limitation of the synthetic environment rather than of the method, and
it is the clearest example in this whole report of what synthetic data cannot
do. A generator can produce data that looks perfectly realistic in its
distributions and still contain none of the structure a model is supposed to
detect. On real subscriber data the ARPU term would be expected to carry most of
the signal. Here it carries none, and no amount of tuning would ever reveal a
pattern that was never generated in the first place.

Which means, being blunt about it, that most of the weighting is doing nothing
here. ARPU decline at 0.35, data decline at 0.15 and voice decline at 0.10 come
to 0.60 of the total, and all three rest on the same missing structure. What is
actually separating the bands is unpaid bills, engagement and tenure --- the
remaining 0.40. On real data those decline terms would be expected to carry most
of the signal, which is why they are weighted the way they are; in this world
they carry none, and no amount of tuning would reveal a pattern that was never
generated in the first place.

It is worth being equally clear about what this component is and is not. It
ranks subscribers by risk factors that can be observed and explained, which is
useful for prioritising who to look at first. It is not a validated predictor of
churn, and Chapter~\ref{ch:limits} sets out why it cannot be one in this
environment. It is one data source among several that the agent can reach and
reason over --- which is the actual subject of this report --- rather than a
result in its own right.


\section{Threats to Validity}

Three limits on how far these numbers should be read.

\textbf{The golden set is small.} Sixteen questions cannot characterise the space
of questions an operator might ask. It is enough to establish that factual
retrieval and safety hold and that hallucination resistance does not; it is not
enough to state a confidence interval on any of those.

\textbf{The environment is synthetic.} Ground truth is computable precisely
because the data was generated. On a real network, the reference itself would be
contested, and the same harness would be considerably harder to build.

\textbf{The agent is not deterministic.} The same question can take a different
query path on a different run. During this evaluation one question was observed
answering correctly on one run and incorrectly on another, with no change in
between. A single pass therefore measures one sample of the agent's behaviour, not
its behaviour.
